%% kapa: exam=1 round=34 year=2023 period=1 form=A subject=real_estate subject_ko=부동산학원론 seq=25 number=105
%% source: Q-Net 감정평가사 기출문제 (cst00309) artlSeq=5211789 file=raw/1008/5211789/2023년도 제34회 감정평가사 1차시험 1교시.pdf
%% conversion: kordoc 4.24.51 -> tools/gen1.py (auto); review=proofread-sonnet; flags=none
\begin{question}{105}
다음 설명에 모두 해당하는 부동산개발방식은?
\begin{framedtext}
○ 사업부지를 소유하고 있는 토지소유자가 개발이 완료된 후 개발업자나 시공사에게 공사대금을 완공된 일부의 건물로 변제하고, 나머지는 분양하거나 소유하는 형태이다. \par
○ 토지소유자는 대상 부지의 소유권을 소유한 상태에서 개발사업이 진행되도록 유도할 수 있고, 그 결과 발생되는 부동산가치의 상승분을 취득할 수 있는 이점이 있다. \par
\end{framedtext}
\begin{choices}
  \item 공영개발방식
  \item 직접개발방식
  \item 대물교환방식
  \item 토지신탁방식
  \item BTL사업방식
\end{choices}
\end{question}
\answer{3}
